\ifdefined\CRevisionRound\else\def\CRevisionRound{2}\fi
\documentclass[10pt]{article}
\pdfvariable suppressoptionalinfo 767
\usepackage[a4paper,margin=20mm]{geometry}
\usepackage{amsmath,amssymb,amsthm,booktabs}
\usepackage{fontspec,unicode-math}
\setmainfont{TeX Gyre Pagella}
\setmathfont{TeX Gyre Pagella Math}
\setmonofont{Latin Modern Mono}
\newfontfamily\cjkfont{Droid Sans Fallback}
\usepackage[hidelinks]{hyperref}
\newcommand{\ii}{\mathrm i}
\newcommand{\e}{\mathrm e}
\newcommand{\dd}{\,\mathrm d}
\newcommand{\R}{\mathbb R}
\newcommand{\C}{\mathbb C}
\newcommand{\Z}{\mathbb Z}
\DeclareMathOperator{\Ree}{Re}
\DeclareMathOperator{\dom}{Dom}
\newtheorem{theorem}{Theorem}
\newtheorem{proposition}{Proposition}
\newtheorem{lemma}{Lemma}
\ifcase\CRevisionRound
\title{Supercritical Inverse-Square Dynamics:\\Domains, Collisions and the Complete Bound Ladder}
\or
\title{Supercritical Inverse-Square Dynamics:\\Complete Spectral Resolution and Scattering}
\else
\title{Supercritical Inverse-Square Dynamics:\\Domain-Scale Cycles and the Ordinary-Determinant Obstruction}
\fi
\author{HCS-C391 theorem and reproducibility package}
\date{5 September 2026}
\begin{document}
\emergencystretch=2em
\maketitle
\begin{abstract}
\ifcase\CRevisionRound
We reconstruct the complete self-adjoint boundary family of the strictly
supercritical inverse-square Hamiltonian on the positive half-line.
Its classical counterpart has an exact quadratic position law, and every
trajectory reaches an excluded collision endpoint in at least one time direction.
Endpoint analysis gives a circle of quantum boundary conditions, none selected
by the classical differential expression. The Bessel Green function determines
every simple negative level and its normalized eigenfunction.
The levels form a bilateral geometric sequence with ratio
$\exp(-2\pi/\sigma)$, accumulating at zero in one direction and diverging to
minus infinity in the other. In particular, no self-adjoint realization has a
ground state or a Friedrichs construction. The normalization is derived from
a two-momentum Lagrange identity, and branch changes only relabel the levels.
The formulas retain their classical literature ownership. This first draft
closes the source clock, all boundary domains and the complete bound ladder;
finite rational and special-function checks are distinguished from proof.
\or
We give a convention-complete spectral reconstruction of every unit-modulus
boundary realization of the supercritical half-line inverse-square Hamiltonian.
The classical collision flow and its quantum self-adjoint extension remain
distinct sources of time evolution. An explicit Green function yields the
entire bilateral negative ladder, while its positive boundary jump gives a
continuous spectrum of multiplicity one and excludes hidden singular spectrum.
Incoming waves are normalized to $(2\pi)^{-1/2}\exp(-\ii kx)$.
The resulting reflection amplitude has unit modulus, but differs by a minus
sign from scattering relative to the free Dirichlet half-line.
The hypotheses of the classical complete-wave-operator theorem are checked
for every positive coupling parameter and every unit boundary phase.
Independent high-precision tests compare 108 Green-kernel jumps with the
spectral-density formula. These tests audit constants and phases; completeness
comes from endpoint and spectral analysis. No target arithmetic interpretation
or novel ownership of the classical inverse-square spectrum is claimed.
\else
We assemble a complete, domain-sensitive theorem for the strictly supercritical
inverse-square Hamiltonian on the positive half-line. Every unit boundary phase
defines a self-adjoint realization with a simple bilateral geometric negative
ladder and a multiplicity-one absolutely continuous positive spectrum.
Its normalized scattering amplitude is explicit. Dilation moves the boundary
phase around a circle; one fixed self-adjoint domain is preserved exactly by
the discrete subgroup $(\pi/\sigma)\Z$, and scattering has the corresponding
least logarithmic momentum period. This domain cycle is not a physical-time
return. The classical source reaches an excluded collision point, whereas
the quantum group is complete and has no nonzero full return time.
The negative ladder makes every positive-time heat operator unbounded;
the continuum makes every resolvent noncompact; and the two ladder tails
leave no convergence germ for an ordinary bilateral spectral zeta.
Thus natural self-adjoint quantization alone supplies none of these global
determinant constructions. The classical formulas are attributed, and exact
algebra, independent Green/Stone checks and adversarial release tests delimit
the package's reproducible contribution.
\fi
\end{abstract}
\noindent\textbf{Keywords:} inverse-square Hamiltonian; self-adjoint domain;
bilateral bound spectrum; continuous scattering; discrete scaling; determinant obstruction.

\medskip
\noindent{\cjkfont\small\raggedright 中文摘要：\par
\ifcase\CRevisionRound
本文完整整理半直线超临界反平方势的自伴边界族。经典位置平方满足精确二次律，
每条轨道至少在一个时间方向抵达被排除的碰撞端点。量子边界构成一个圆，
并不由经典方程唯一选定。贝塞尔核给出全部简单负本征值及归一化函数。
双向几何能级分别趋于零和负无穷，因此没有基态或弗里德里希斯实现。
本文保留经典文献归属，有限核验不替代完整证明。
\or
本文在全部自伴边界与负能级结论之外，加入完整连续谱分解和散射规范。
格林核的边界跳跃确定单重绝对连续谱；入射波归一化明确区分反射振幅与相对散射。
独立高精度核验直接比较谱密度与格林核跳跃，防止常数和相位错误。
完备性来自算子证明及已核实条件的文献定理，而非有限计算。
\else
本文将全部自伴边界、双向负能级、连续谱散射与离散尺度循环闭合为一个定理包。
尺度变换改变边界相位，固定定义域只在离散子群下保持不变，但这不是物理时间周期。
负谱无下界导致热演化无界，连续谱阻止预解算子紧性，双向能级级数没有收敛起点。
因此自然自伴量子化不能直接产生普通全局行列式，更不等于目标算术谱的实现。
经典公式保留原作者归属，计算与交叉审查仅承担明确的复算核验任务。
\fi
\par\noindent 关键词：反平方势；自伴定义域；双向负谱；连续散射；离散尺度；行列式障碍\par}

\section{A local Hamiltonian does not choose a boundary}
Fix $\sigma>0$, put $g=\sigma^2+1/4$, and use units $\hbar=2m=1$.
The classical source is $h(x,p)=p^2-g/x^2$ on $T^*(0,\infty)$.
Its quantum expression is $\ell=-\partial_x^2-g/x^2$ on
$\mathcal H=L^2((0,\infty),\dd x)$. A unit complex number $\kappa$ will
specify the missing boundary condition. It is independent physical input,
not a value fitted to a desired eigenvalue or prime scale.
Throughout, $I_\nu,J_\nu,K_\nu$ are ordinary Bessel functions.

The contribution is a single-convention source theorem and audit, not a new
discovery of the inverse-square spectrum. Dereziński and Richard give the
domain, spectrum and scattering theory in a substantially larger complex
parameter family \cite{dr2017}. We retain the entire strictly supercritical
self-adjoint subfamily, derive the constants used here, and connect them to
the original source clock and the limits of ordinary global determinants.
The critical value $\sigma=0$, nonunit $\kappa$, ultraviolet regularization
and confinement are not part of this theorem.

\begin{proposition}[The complete collision clock]
For every initial point $(x_0,p_0)$, with energy $E=h(x_0,p_0)$,
\begin{equation}
 y(t)=x(t)^2=x_0^2+4x_0p_0t+4Et^2,\qquad
 p(t)=\frac{y'(t)}{4\sqrt{y(t)}}.
 \label{eq:classical}
\end{equation}
Its maximal time interval is the component of $\{t:y(t)>0\}$ containing zero.
Every orbit has a finite collision endpoint, and none is periodic in the
declared phase space.
\end{proposition}
\begin{proof}
Hamilton's equations are $x'=2p$ and $p'=-2g/x^3$, so $y''=8E$.
The reconstruction in \eqref{eq:classical} satisfies both equations.
For $E\ne0$ the quadratic discriminant is $16g>0$.
At $E<0$ its positive component is bounded; at $E>0$ it is a half-line.
At $E=0$ the nonzero slope is $\pm4\sqrt g$, again giving a half-line.
Each finite endpoint has $x=0$, which is excluded. No collision continuation
has been added, so none of these intervals gives a periodic orbit.
\end{proof}

\section{All self-adjoint domains and the bound ladder}
Let $D_{\max}=\{f\in L^2:\ell f\in L^2\text{ distributionally}\}$ and
let $D_{\min}$ be the graph closure of $C_c^\infty(0,\infty)$.
Membership in $D_{\min}$ near zero means membership after multiplication
by a smooth cutoff equal to one there. Define
\begin{equation}
 D_\kappa=\left\{f\in D_{\max}:\ %
 f-c(\kappa x^{1/2-\ii\sigma}+x^{1/2+\ii\sigma})
 \in D_{\min}\text{ near }0\text{ for some }c\in\C\right\}.
 \label{eq:domain}
\end{equation}

\begin{theorem}[Exhaustive boundary family]
The domains \eqref{eq:domain}, $|\kappa|=1$, are exactly the self-adjoint
extensions of the minimal operator. None is semibounded, and there is no
Friedrichs realization. Position-space complex conjugation preserves every
one of these domains.
\end{theorem}
\begin{proof}
The zero-energy basis $u_\pm=x^{1/2\pm\ii\sigma}$ is square integrable
near zero. Variation of constants gives a unique expansion
$f=a_fu_-+b_fu_+$ modulo $D_{\min}$ there.
Indeed, the integrals of $u_\pm\ell f$ converge by Cauchy--Schwarz and
their remainders are $o(x^{3/2})$, with derivatives $o(x^{1/2})$.
Infinity is limit point since the potential is integrable on tails.
Weyl's endpoint theorem therefore gives deficiency indices $(1,1)$.
The boundary Wronskian is
\[
 W_0(\overline f,v)=2\ii\sigma
 (\overline{b_f}b_v-\overline{a_f}a_v).
\]
Its maximal isotropic lines are exactly $a=\kappa b$, $|\kappa|=1$;
neither coefficient of a nonzero isotropic vector can vanish.
Conjugation changes the ratio to $1/\overline\kappa=\kappa$.
For nonsemiboundedness, set $f(x)=x^{1/2}\chi(\log x)$ with a real compactly
supported smooth function $\chi$. Direct substitution gives
\[
 \langle f,\ell f\rangle=\int_{\R}
 (|\chi'(s)|^2-\sigma^2|\chi(s)|^2)\dd s.
\]
A sufficiently long plateau makes this negative. Normalized dilations make
it arbitrarily negative. Every extension contains these minimal-domain
test functions, so no semibounded extension exists.
\end{proof}

Put $\varsigma=\kappa\Gamma(-\ii\sigma)/\Gamma(\ii\sigma)=\e^{\ii\theta}$,
using any real lift $\theta$. Powers of a right-half-plane momentum use
the logarithm on $\Ree\rho>0$.
\begin{theorem}[The entire negative spectrum]
Every negative spectral point of $H_{\sigma,\kappa}=\ell|_{D_\kappa}$
is a simple eigenvalue
\begin{equation}
 E_j=-\rho_j^2=-4\exp\!\left(-\frac{\theta+2\pi j}{\sigma}\right),
 \quad j\in\Z,\qquad
 \psi_j(x)=\rho_j\sqrt{\frac{2\sinh(\pi\sigma)}{\pi\sigma}}
 \sqrt x K_{\ii\sigma}(\rho_jx).
 \label{eq:ladder}
\end{equation}
The functions $\psi_j$ have unit norm. A different lift of $\theta$ merely
reindexes the same spectrum.
\end{theorem}
\begin{proof}
At energy $-\rho^2$ the decaying solution at infinity is uniquely
$\sqrt xK_{\ii\sigma}(\rho x)$, up to a scalar.
Its endpoint coefficients follow from
\[
 K_{\ii\sigma}(w)=\tfrac12\Gamma(\ii\sigma)(w/2)^{-\ii\sigma}
 +\tfrac12\Gamma(-\ii\sigma)(w/2)^{\ii\sigma}+o(1).
\]
Their ratio equals $\kappa$ exactly when
$\varsigma(\rho/2)^{2\ii\sigma}=1$, giving \eqref{eq:ladder}.
For $\Ree\rho>0$ this equality forces $\arg\rho=0$ by taking absolute values.
Away from these roots the Green kernel in the proof supplement gives a
bounded inverse; hence no additional negative spectrum is omitted.
Uniqueness of the decaying solution gives simplicity.

For $f_a(x)=\sqrt xK_{\ii\sigma}(ax)$ and $f_b(x)=\sqrt xK_{\ii\sigma}(bx)$,
the Lagrange identity is $W(f_a,f_b)'=(b^2-a^2)f_af_b$.
The Wronskian vanishes at infinity; its endpoint value is
$\sigma|\Gamma(\ii\sigma)|^2\sin(\sigma\log(a/b))$.
Using $|\Gamma(\ii\sigma)|^2=\pi/(\sigma\sinh\pi\sigma)$ gives
\begin{equation}
 \int_0^\infty xK_{\ii\sigma}(ax)K_{\ii\sigma}(bx)\dd x
 =\frac{\pi\sin(\sigma\log(a/b))}{(a^2-b^2)\sinh(\pi\sigma)}.
 \label{eq:lagrange}
\end{equation}
The coincident limit is $\pi\sigma/(2a^2\sinh\pi\sigma)$, proving the
normalizer in \eqref{eq:ladder}. Endpoint $O(x)$ bounds and exponential
tail decay justify that limit. Distinct ladder values give orthogonality
in \eqref{eq:lagrange}. Finally $E_j\to0$ as $j\to+\infty$ and
$E_j\to-\infty$ as $j\to-\infty$; there is no lowest level.
\end{proof}

\ifnum\CRevisionRound>0
\section{Green jumps, completeness and scattering}
For $z=-\rho^2$, $\Ree\rho>0$, define
\begin{align}
 T_\rho&=\varsigma(\rho/2)^{2\ii\sigma},\nonumber\\
 u_\rho(x)&=\sqrt x\,[I_{\ii\sigma}(\rho x)-T_\rho I_{-\ii\sigma}(\rho x)],
 &v_\rho(x)&=\sqrt xK_{\ii\sigma}(\rho x).\label{eq:uv}
\end{align}
The Gamma recurrence gives the required endpoint ratio for $u_\rho$.
Both $I$--$K$ Wronskians equal $-1$, so
$W(u_\rho,v_\rho)=T_\rho-1$. Thus
\begin{equation}
 G_\kappa(-\rho^2;x,y)=
 \frac{u_\rho(\min(x,y))v_\rho(\max(x,y))}{1-T_\rho}.
 \label{eq:green}
\end{equation}
Its first derivative jumps by $-1$. The $O(\sqrt{xy})$ behavior on compact
endpoint squares and exponential tail bounds give a bounded inverse when
the denominator is nonzero.

For positive $k$, put $t(k)=\varsigma(k/2)^{2\ii\sigma}$,
$a=\e^{\pi\sigma/2}$ and $b=\e^{-\pi\sigma/2}$. Define
\begin{equation}
 \phi_k(x)=\e^{-\ii\pi/4}\sqrt{kx}\,
 \frac{J_{\ii\sigma}(kx)-t(k)J_{-\ii\sigma}(kx)}{b-t(k)a}.
 \label{eq:phi}
\end{equation}
The denominator satisfies $|b-ta|\ge a-b>0$, uniformly in $k$ at fixed
$\sigma$. The two Green boundary values give
\begin{equation}
 \frac{k}{\pi\ii}\big[G_\kappa(k^2+\ii0;x,y)-G_\kappa(k^2-\ii0;x,y)\big]
 =\phi_k(x)\overline{\phi_k(y)}.
 \label{eq:stone}
\end{equation}

\begin{theorem}[Complete spectral measure]
The positive spectrum is absolutely continuous of multiplicity one, with
momentum density $\phi_k(x)\overline{\phi_k(y)}\dd k$.
There are no positive or zero eigenvalues and no singular continuous spectrum.
For every $f\in\mathcal H$,
\begin{equation}
 \|f\|^2=\sum_{j\in\Z}|\langle\psi_j,f\rangle|^2+
 \int_0^\infty|\langle\phi_k,f\rangle|^2\dd k.
 \label{eq:parseval}
\end{equation}
The transform is initially defined weakly on compactly supported test
functions and extended by this identity.
\end{theorem}
\begin{proof}
The $I$--$J$ and $K$--Hankel identities in \eqref{eq:green} give
\eqref{eq:stone}. On compact positive momentum intervals the boundary values
are continuous between the weighted spaces used in limiting absorption,
with weight exponent greater than $1/2$; endpoint and oscillatory tail
bounds apply because $b-ta$ never vanishes. These are the nonexceptional
hypotheses in \cite[Theorem 6.1 and Proposition 6.7]{dr2017}.
Stone's formula gives the displayed density and rules out singular spectrum
on every such interval. Positive-energy Bessel waves have nondecaying
oscillatory tails, so no nonzero solution is square integrable there.
For every nonzero zero-energy combination, dilation by $\e^{\pi/\sigma}$
multiplies its squared norm on consecutive logarithmic annuli by
$\e^{2\pi/\sigma}>1$, so it is not square integrable at infinity.
A singular measure supported only at zero would
be an eigenvalue atom, which is excluded. The full spectral theorem together
with the negative resolvent classification now yields \eqref{eq:parseval}.
One generalized eigenfunction per momentum gives multiplicity one.
\end{proof}

The large-$x$ expansion, with no unspecified normalization phase, is
\begin{equation}
 \phi_k(x)=\frac{\e^{-\ii kx}+R(k)\e^{\ii kx}}{\sqrt{2\pi}}+O(x^{-1}),
 \qquad R(k)=-\ii\frac{a-t(k)b}{b-t(k)a}.
 \label{eq:reflection}
\end{equation}
Incoming means $\e^{-\ii kx}$. Numerator and denominator in the ratio have
equal modulus, so $|R|=1$. This is reflection, not yet relative scattering.
The free Dirichlet half-line Hamiltonian $H_D$ has reflection $-1$.
Under the convention
\[
 W_\pm=\mathop{\mathrm{s-lim}}_{t\to\pm\infty}
 \e^{\ii tH_{\sigma,\kappa}}\e^{-\ii tH_D},\qquad S=W_+^*W_-,
\]
the sine-representation multiplier is $S(k)=-R(k)$.
Existence and completeness of these limits are the explicitly used external
result \cite[Proposition 6.10]{dr2017}. Its two parameter pairs are
$(\ii\sigma,\kappa)$ and $(1/2,0)$: both real parts lie strictly between
minus one and one, both realizations are self-adjoint, and neither denominator
is exceptional. Proposition 6.9 of that source agrees with the phase in
\eqref{eq:reflection}. Completeness means $W_\pm^*W_\pm=I$ and
$W_\pm W_\pm^*=P_{\rm ac}(H)$, not that negative bound states scatter freely.
\fi

\ifnum\CRevisionRound>1
\section{The scale cycle is not a return of time evolution}
For the unitary dilation $(U_\tau f)(x)=\e^{\tau/2}f(\e^\tau x)$,
substitution into \eqref{eq:domain} gives
\begin{align}
 U_\tau D_\kappa&=D_{\kappa\e^{-2\ii\sigma\tau}},\nonumber\\
 U_\tau H_{\sigma,\kappa}U_\tau^{-1}
 &=\e^{-2\tau}H_{\sigma,\kappa\e^{-2\ii\sigma\tau}}.
 \label{eq:scale}
\end{align}
Hence a fixed self-adjoint domain is invariant exactly when
$\tau\in(\pi/\sigma)\Z$. No self-adjoint domain is continuously homogeneous.
This cycle is in domain space; it does not use the original Hamiltonian time.
The scattering formula has
\[
 t(k\e^{\pi/\sigma})=t(k),\qquad
 R(k\e^{\pi/\sigma})=R(k),\qquad S(k\e^{\pi/\sigma})=S(k).
\]
The logarithmic period $\pi/\sigma$ is the least positive one, because the
Möbius map $t\mapsto(a-tb)/(b-ta)$ is injective:
its determinant is $a^2-b^2\ne0$. The same scale produces the bound ratio
$|E_{j+1}|/|E_j|=\e^{-2\pi/\sigma}$. Neither scale has a rational-prime origin.

\begin{theorem}[The ordinary global-determinant obstruction]
For every unit boundary parameter, the group $\e^{-\ii tH}$ exists for every
real time and has no nonzero full return time. For every $t>0$ the operator
$\e^{-tH}$ is unbounded. Every resolvent of $H$ is noncompact and therefore
belongs to no Schatten class. The bilateral series
$\sum_{j\in\Z}(-E_j)^{-s}$ converges at no complex value of $s$.
\end{theorem}
\begin{proof}
Self-adjointness gives the unitary group. On the nonzero absolutely continuous
subspace, multiplication by $\e^{-\ii tk^2}$ cannot be identically one unless
$t=0$. The complete group is thus distinct from both the collision flow and
the domain dilation cycle. Since $E_j\to-\infty$, the heat values
$\e^{-tE_j}$ diverge along normalized eigenvectors for every $t>0$.

On a bounded positive spectral interval the resolvent is multiplication
by $(k^2-z)^{-1}$, bounded away from zero. Infinitely many normalized
functions with disjoint spectral supports have images separated in norm;
therefore that resolvent is not compact. No ordinary Fredholm determinant
of $I-w(H-z)^{-1}$ is available.

For the bilateral series, terms along $j\to+\infty$ fail to vanish when
$\Ree s\ge0$ and terms along $j\to-\infty$ fail when $\Ree s\le0$.
Both conditions are necessary, so no value of $s$ works. The corresponding
unregularized product $\prod_j(1-w/E_j)$ has no nonzero neighborhood of
convergence, since its factors fail to tend to one along $E_j\to0$.
\end{proof}

\begin{samepage}
The obstruction concerns \emph{ordinary} constructions. It does not prohibit
relative resolvents or explicitly renormalized determinants, and it does not
identify any such different object with an arithmetic target.
The strict tuple is
\[
 (\mathrm{A0\_FAIL},\mathrm{A1\_FAIL},\mathrm{A2\_FAIL},
   \mathrm{A3\_FAIL},\mathrm{A4\_NATURAL\_QUANTIZATION}),
\]
with overall \texttt{ROUTE\_A\_REJECTED}.
A4 records only the natural differential expression and the proved domains;
the free boundary parameter is not selected by the classical source.
No arithmetic local data, target Euler factor, root number, automorphy,
target divisor or counting law, target functional equation, target-zero match,
or Hilbert--Pólya operator is claimed. Route B remains disabled.
The scope is \texttt{NO\_BAD\_EULER\_OR\_ROOT\_NUMBER}.
\end{samepage}
\fi

\section{Reproducibility and literature ownership}
The exact producer contains 45 classical phase points, 27 boundary-flux rows
and 15 algebraic scattering rows. A non-importing SymPy checker reconstructs
their values and literal scalar types. Sixty bound-level and 36 continuum
rows are stored to 60 digits after 100-digit evaluation and independently
recomputed at 110 digits through endpoint-normalized coefficients.
The separate symbolic lane checks eight identities.
\ifnum\CRevisionRound>0
The 100-digit special-function lane checks 108 actual Stone-jump cells,
12 Green Wronskians, 12 differential-equation residuals, 12 endpoint matches,
three bound-state normalization integrals and 36 logarithmic-period cells.
In particular, testing $|R|=1$ alone is not treated as a spectral-density test.
\fi
These finite grids are regression, not interval certification or a proof of
spectral completeness. Two unrelated working-directory replays, repaired-hash
attacks, strict YAML gates and smoke tests are distributed with the proof.
The release rebuilds all three drafts in fresh directories with frozen epoch
1788566400 and retains actual settled compiler logs.

The external mathematical ownership is concentrated in one directly
inspected primary source \cite{dr2017}; listing its journal and preprint as
two independent works would be misleading. The paper does not claim priority
for the self-adjoint family, geometric ladder or scattering formulas.
Its package-level advance is the joined, auditable boundary between that
complete source theorem and a target spectral construction it does not provide.
The neighboring repository models are a Friedrichs Aharonov--Bohm Hamiltonian,
a repulsive many-body Calogero--Moser system and a positive isotonic oscillator.
They do not contain this strictly supercritical, imaginary-order, all-domain
bilateral ladder. The proof and source audit preserve those owner differences.

Local collaborating AI agents reviewed the proof and its normalization;
this is not external human peer review or cross-provider validation.
An independent reading prompted direct Green/Stone tests and a fully expanded
normalization identity, both implemented here. No outside model was invoked.

\section{Conclusion}
\ifcase\CRevisionRound
Round zero: domains and the complete bound ladder.
The original source has an explicit collision clock and a complete circle
of quantum self-adjoint domains. The bilateral negative spectrum has no
lowest level. This closes the endpoint and bound-state foundation without
pretending that a classical collision law selected a boundary phase.
\or
Round one: complete continuum and scattering.
The Green boundary jump and endpoint analysis close both spectral components.
The normalized reflection and relative scattering phases remain distinct,
and the precise external wave-operator theorem supplies time-dependent
completeness only on the absolutely continuous subspace.
\else
Round two: scale cycles and ordinary-determinant obstruction.
The domain circle, bilateral ladder and continuous scattering fit one
convention-complete theorem. Natural self-adjointness gives a unitary group,
but neither a ground state, a bounded heat operator, compact resolvent nor
an ordinary bilateral-zeta germ. The negative Route-A verdict is part of the
result, not a reason to change the source clock or insert target data.
\fi

\begin{thebibliography}{1}
\raggedright
\bibitem{dr2017}
J.~Dereziński and S.~Richard.
On Schrödinger operators with inverse square potentials on the half-line.
\emph{Annales Henri Poincaré} \textbf{18}, 869--928 (2017).
\href{https://doi.org/10.1007/s00023-016-0520-7}{doi:10.1007/s00023-016-0520-7}.
Author manuscript: \href{https://arxiv.org/abs/1604.03340v2}{arXiv:1604.03340v2}.
The theorem numbering used here follows that inspected preprint.
\end{thebibliography}
\end{document}
